\chapter[Algorithm-Level Study of Nine ML Methods]{Algorithm-Level Study of\\ Nine Machine-Learning Methods}\label{ch:algorithms}

Section~7.4 of the project proposal named nine candidate algorithms: linear regression, logistic regression,
decision tree, random forest, gradient boosting, XGBoost, support vector machine, $k$-nearest neighbours and
$K$-means clustering. \Cref{ch:models} used them inside the production pipeline, with rich feature sets,
randomised hyper-parameter search and statistical benchmarks. This chapter takes a complementary view. Each
algorithm is studied \emph{on its own}, in a separate and fully transparent experiment, in order to show
\emph{how} it learns, which settings control it, how it fails and what it contributes to the business problem.
The experiments also serve as an independent check of \cref{ch:models}: they use a deliberately simpler feature
set and separate code, so agreement between the two chapters strengthens the conclusions.

\section{Common experimental protocol}\label{sec:algoprotocol}
\subsection{Three problems}
All nine experiments use the same data and the same rules, implemented once in
\path{algorithms/common.py}:
\begin{itemize}
  \item \textbf{Problem A -- regression.} Monthly units of each product. Transactions are aggregated into one row
        per product and month, described by nine features that are known before the month starts: month, a
        December indicator, product age, mean log price, mean discount, pack weight, and three history features
        (units last month, units in the same month last year, and the mean of the previous three months).
  \item \textbf{Problem B -- classification.} Whether an order is large ($\geq 10$ units), predicted from the 18
        order-level features selected in the feature-engineering step (calendar, price, pack, history and the
        smoothed channel~$\times$~region order size). Large orders form \LargeShareTrain of training orders and
        \LargeShareTest of test orders, so the classes are imbalanced. A fixed random sample of 15,000 orders
        from 2020--2024 and 15,000 from 2025--2026 is used by every algorithm. The sample keeps the slower
        algorithms (SVM, KNN) practical and makes the comparison fair.
  \item \textbf{Problem C -- clustering.} One profile per product (\KMProducts products) from 2020--2024: average
        price, units per month on sale, average order size, average discount, share of large orders and the
        share of December sales. There is no target variable.
\end{itemize}

\subsection{Four-way time split}
Every supervised experiment follows the same sequence (\cref{fig:split}). Settings are chosen on 2024
(validation) after learning from 2020--2023. The chosen setting is then refitted on 2020--2024, and the test
period, January 2025 -- September 2026, is scored once. The test data therefore never influence any choice.
Every model is compared with a naive baseline: ``units this month = units last month'' for problem~A and
``always small'' for problem~B.

\begin{figure}[htbp]\centering
\begin{tikzpicture}[x=1.75cm, every node/.style={font=\footnotesize}]
\foreach \y/\lab in {0/2020,1/2021,2/2022,3/2023,4/2024,5/2025,6/2026} \node at (\y+0.5,0.85) {\lab};
\fill[brand!25] (0,0) rectangle (4,0.5);   \node at (2,0.25) {\textbf{fit} -- learn while tuning};
\fill[accent!35] (4,0) rectangle (5,0.5);  \node at (4.5,0.25) {\textbf{val}};
\fill[brand!45] (0,-0.75) rectangle (5,-0.25); \node at (2.5,-0.5) {\textbf{train} -- learn again with the chosen setting};
\fill[green!30!gray!40] (5,-0.75) rectangle (6.75,-0.25); \node at (5.875,-0.5) {\textbf{test}};
\draw[thick,-{Stealth}] (0,-1.05) -- (7,-1.05) node[right] {time};
\node[align=left, anchor=west, font=\scriptsize] at (0,-1.45) {1. choose settings on 2024 \quad 2. refit on 2020--2024 \quad 3. score once on 2025--2026};
\end{tikzpicture}
\caption{Time-ordered protocol used by every supervised experiment in this chapter}\label{fig:split}
\end{figure}

\subsection{Implementation}
Each algorithm has its own script in the folder \path{algorithms/} (\cref{tab:lessons}). Every script is divided
into small numbered steps, so it can be run as a whole or cell by cell. All steps print their results and save
their charts to \path{algorithms/outputs/}. The final test scores and the key numbers quoted below are written
to JSON files, and \path{algorithms/export_to_report.py} turns them into the tables and \LaTeX{} macros of this
chapter. As in the rest of the report, no number here is typed by hand.

\begin{table}[htbp]\centering\small
\caption{Algorithm experiments and what each one demonstrates}\label{tab:lessons}
\begin{tabularx}{\linewidth}{@{}llX@{}}
\toprule
\textbf{Script} & \textbf{Problem} & \textbf{Concepts demonstrated} \\ \midrule
\texttt{01\_linear\_regression} & A & learned equation, standardised coefficients, residual analysis \\
\texttt{02\_logistic\_regression} & B & sigmoid, odds ratios, class weights, threshold choice, ROC curve \\
\texttt{03\_decision\_tree} & B (+A) & readable rules, Gini impurity, depth and overfitting \\
\texttt{04\_random\_forest} & A (+B) & bagging, number of trees, out-of-bag score, permutation importance \\
\texttt{05\_gradient\_boosting} & A (+B) & boosting by hand, learning rate versus number of trees \\
\texttt{06\_xgboost} & B (+A) & regularisation, early stopping, rare-class weighting, gain importance \\
\texttt{07\_svm} & B (+A) & scaling, linear and RBF kernels, the $C$ and $\gamma$ settings, SVR \\
\texttt{08\_knn} & B (+A) & distance, scaling, choice of $k$, inspecting neighbours \\
\texttt{09\_kmeans\_clustering} & C & Lloyd's algorithm by hand, elbow and silhouette, naming clusters \\ \bottomrule
\end{tabularx}
\end{table}

% ============================================================================
\section{Linear regression}
\paragraph{Principle.} Linear regression models the target as a weighted sum of the inputs,
$\hat y = \beta_0 + \beta_1 x_1 + \dots + \beta_p x_p$. Ordinary least squares chooses the weights that minimise
the sum of squared errors, which has the closed-form solution
\[
  \hat{\boldsymbol\beta} = (\mathbf X^\top \mathbf X)^{-1}\mathbf X^\top \mathbf y .
\]
Each coefficient is the expected change in the target for a one-unit change in its input, holding the others
constant \citep{montgomery2021, james2021}.

\paragraph{Experiment.} With last month's units as the only input, the learned line is
$\widehat{\text{units}} = \LinBzero + \LinBone \times \text{units}_{t-1}$ (\cref{fig:alglin}, left). The slope
below one is \emph{regression to the mean}: an unusually strong month tends to be followed by a weaker one, and
vice versa. This single-input model already reduces the test MAE from \AlgRegBaseMAE (baseline) to \LinOneMAE
units. With all nine standardised inputs, the test MAE falls to \AlgLinMAE units ($R^2 = \AlgLinRtwo$). The
largest standardised coefficient belongs to \texttt{\LinTopFeature} (\LinTopCoef units per standard
deviation): the recent level of demand is the strongest predictor. The residual plot (\cref{fig:alglin}, right)
is centred on zero with no strong pattern, so the linear form is adequate for this problem.

\begin{figure}[htbp]\centering
\includegraphics[width=.43\linewidth]{algo/01_simple_line}\hfill
\includegraphics[width=.55\linewidth]{algo/01_residuals}
\caption{Linear regression: the learned line with one input (left); predicted versus actual units and the error distribution on the test period (right)}\label{fig:alglin}
\end{figure}

\paragraph{Assessment.} Linear regression is fast, fully explainable and, with good lag features, nearly as
accurate as the ensembles. Its weaknesses are that it cannot represent interactions or thresholds unless they
are added as features, and that correlated inputs (here the three history features) make individual
coefficients unstable even when predictions are good.

% ============================================================================
\section{Logistic regression}
\paragraph{Principle.} Logistic regression passes the linear score $z = \beta_0 + \boldsymbol\beta^\top\mathbf x$
through the sigmoid function to obtain a probability,
\[
  P(y=1 \mid \mathbf x) = \sigma(z) = \frac{1}{1+e^{-z}}, \qquad
  \log\frac{P}{1-P} = \beta_0 + \boldsymbol\beta^\top\mathbf x ,
\]
and the weights maximise the log-likelihood of the observed classes \citep{hosmer2013}. The exponentiated
coefficient $e^{\beta_j}$ is an \emph{odds ratio}: the factor by which the odds of a large order are multiplied
when $x_j$ increases by one unit.

\paragraph{Experiment.} With only the smoothed channel~$\times$~region order size as input, the coefficient is
\LogBone. The predicted probability of a large order rises from \LogPfive in a channel and region whose orders
average 5 units, to \LogPten at 10 units and \LogPfifteen at 15 units. With all 18 standardised inputs and
\emph{balanced class weights} (errors on the rare class count more), the same feature dominates with an odds
ratio of \LogTopOR per standard deviation. All other odds ratios are close to one. The decision threshold was
chosen by maximising F1 on 2024, which selected \LogThreshold. On the test period the model reaches ROC-AUC
\AlgLogAUC, recall \AlgLogRecall and precision \AlgLogPrec. The baseline that always answers ``small'' is
\AlgClsBaseAcc\% accurate but finds no large order at all. This shows why accuracy alone is misleading for
imbalanced classes.

\begin{figure}[htbp]\centering
\includegraphics[width=.36\linewidth]{algo/02_sigmoid}\hfill
\includegraphics[width=.62\linewidth]{algo/02_confusion_roc}
\caption{Logistic regression: the sigmoid link (left); confusion matrix and ROC curve on the test period (right)}\label{fig:alglog}
\end{figure}

\paragraph{Assessment.} Logistic regression gives calibrated probabilities and explanations that managers can
read directly (``this channel doubles the odds''). In this study it is also the most accurate classifier, which
is strong evidence that the signal is essentially one main effect.

% ============================================================================
\section{Decision tree}
\paragraph{Principle.} A decision tree splits the data recursively with yes/no questions on single features. At
each node, CART chooses the question that most reduces impurity \citep{breiman1984}. For classification the
Gini impurity of a node with class shares $p_k$ is $G = 1 - \sum_k p_k^2$, which is zero for a pure node. For
regression the variance (squared error) is used, and each leaf predicts the mean of its training rows.

\paragraph{Experiment.} A tree of depth two already reads like a business rule: \emph{if the channel and region
usually place orders above about six units, expect a large order}. The depth was tuned on 2024
(\cref{fig:algdt}, left). As the tree deepens, its ROC-AUC on the data it learned from keeps rising (to
\DTFitAUCDeep at depth 15), while the score on unseen 2024 orders peaks early and then falls (\DTValAUCDeep at
depth 15). This is the textbook picture of \emph{overfitting}. The best depth was \DTBestDepth, with a test
ROC-AUC of \AlgDTAUC. In the final tree \texttt{channel\_region\_avg\_qty} accounts for \DTTopShare of the
total impurity reduction. As a regressor for problem~A, a depth-5 tree reaches an MAE of \AlgDTMAE units: better
than the baseline, but the weakest of the supervised regressors.

\begin{figure}[htbp]\centering
\includegraphics[width=.42\linewidth]{algo/03_depth_tuning}\hfill
\includegraphics[width=.56\linewidth]{algo/03_tree}
\caption{Decision tree: tuning the depth on 2024 (left); the top levels of the final tree (right)}\label{fig:algdt}
\end{figure}

\paragraph{Assessment.} Trees are the most explainable non-linear model and need no scaling, but a single tree has
high variance: small changes in the data can produce a different tree. The two ensemble methods that follow
are designed to remove this weakness.

% ============================================================================
\section{Random forest}
\paragraph{Principle.} A random forest averages $B$ trees, each grown on a bootstrap sample of the rows and
choosing every split from a random subset of the features \citep{breiman2001}. If single trees have variance
$\sigma^2$ and pairwise correlation $\rho$, the variance of their average is
\[
  \operatorname{Var}\Bigl(\tfrac{1}{B}\textstyle\sum_{b=1}^{B} T_b\Bigr) = \rho\,\sigma^2 + \frac{1-\rho}{B}\,\sigma^2 ,
\]
so adding trees removes the second term, and random feature selection reduces $\rho$ \citep{hastie2009}.

\paragraph{Experiment.} A single full-depth tree has a validation MAE of \RFOneTreeMAE units on 2024. A forest
of one tree is similar (\RFValMAEone), but the error falls quickly as trees are added and levels off at about
100 trees (\RFValMAEhundred; \cref{fig:algrf}). More trees never cause overfitting; they only cost time. Each
tree leaves out about a third of the rows, and scoring each row with the trees that did not see it gives a free
\emph{out-of-bag} estimate ($R^2 = \RFOOB$). On the test period the 300-tree forest reaches an MAE of
\AlgRFMAE units ($R^2 = \AlgRFRtwo$). Permutation importance shows that shuffling
\texttt{\RFTopFeature} increases the error by \RFTopPerm units, more than any other input. The individual trees
also give a range for each forecast, which is useful for safety stock. As a classifier, the forest reaches
ROC-AUC \AlgRFAUC.

\begin{figure}[htbp]\centering
\includegraphics[width=.6\linewidth]{algo/04_n_trees}
\caption{Random forest: validation error (2024) against the number of trees}\label{fig:algrf}
\end{figure}

% ============================================================================
\section{Gradient boosting}
\paragraph{Principle.} Boosting builds trees \emph{in sequence}. Starting from a constant $F_0$, each new tree
$h_m$ is fitted to the negative gradient of the loss, which for squared error is simply the current residual,
and a fraction $\nu$ (the learning rate) of it is added \citep{friedman2001}:
\[
  F_m(\mathbf x) = F_{m-1}(\mathbf x) + \nu\, h_m(\mathbf x), \qquad
  h_m \approx \arg\min_h \textstyle\sum_i \bigl(y_i - F_{m-1}(\mathbf x_i) - h(\mathbf x_i)\bigr)^2 .
\]

\paragraph{Experiment.} The experiment first implements boosting by hand in six lines of code. Starting from the
mean, the training MAE falls from \GBManualStart to \GBManualFive units after only five depth-3 trees. The
library version was then tuned over the learning rate and the number of trees using staged predictions on 2024
(\cref{fig:alggb}). A high learning rate reaches its best error quickly and then overfits. A low rate needs many
more trees but degrades more slowly. The chosen setting was a learning rate of \GBRate with \GBTrees trees. The
test MAE is \AlgGBMAE units ($R^2 = \AlgGBRtwo$), and the most important input is
\texttt{\GBTopFeature}. Because the classifier version has no class-weight option, its decision threshold was
lowered to the training share of large orders. It then reaches ROC-AUC \AlgGBAUC.

\begin{figure}[htbp]\centering
\includegraphics[width=.65\linewidth]{algo/05_learning_rate}
\caption{Gradient boosting: validation error (2024) against the number of trees for three learning rates}\label{fig:alggb}
\end{figure}

% ============================================================================
\section{XGBoost}
\paragraph{Principle.} XGBoost is a regularised, second-order implementation of gradient boosting
\citep{chen2016}. At step $m$ it minimises
\[
  \mathcal L^{(m)} = \sum_i \ell\bigl(y_i, \hat y_i^{(m-1)} + f_m(\mathbf x_i)\bigr) + \Omega(f_m), \qquad
  \Omega(f) = \gamma T + \tfrac12 \lambda \lVert \mathbf w\rVert^2 ,
\]
where $T$ is the number of leaves and $\mathbf w$ the leaf values. The penalty $\Omega$ discourages complex trees.
Row and column sub-sampling add randomness, as in a random forest, and \emph{early stopping} ends training as
soon as the validation score stops improving.

\paragraph{Experiment.} For problem B, the weight of the rare class was set to the ratio of small to large orders
(\XGBPosWeight). With early stopping monitored on 2024, the validation ROC-AUC peaked after only \XGBBestN trees
(\XGBValAUC; \cref{fig:algxgb}). The training curve kept rising after that point, the familiar sign of
memorisation. Refitted with that many trees, XGBoost reaches a test ROC-AUC of \AlgXGBAUC and an F1 of
\AlgXGBFone. Gain importance attributes \XGBTopGain of the model's gain to the channel~$\times$~region feature.
For problem~A early stopping selected \XGBRegTrees trees, and the regressor reaches a test MAE of \AlgXGBMAE
units ($R^2 = \AlgXGBRtwo$).

\begin{figure}[htbp]\centering
\includegraphics[width=.65\linewidth]{algo/06_early_stopping}
\caption{XGBoost: ROC-AUC on the training and validation years as trees are added, with the early-stopping point}\label{fig:algxgb}
\end{figure}

% ============================================================================
\section{Support vector machine}
\paragraph{Principle.} A support vector machine finds the boundary with the widest margin between the classes.
The soft-margin formulation allows some violations, controlled by $C$ \citep{cortes1995}:
\[
  \min_{\mathbf w, b, \boldsymbol\xi}\ \tfrac12 \lVert\mathbf w\rVert^2 + C\sum_i \xi_i
  \quad\text{s.t.}\quad y_i(\mathbf w^\top\phi(\mathbf x_i) + b) \geq 1 - \xi_i,\ \ \xi_i \geq 0 .
\]
The kernel trick replaces inner products by a kernel such as the radial basis function
$K(\mathbf x, \mathbf x') = \exp(-\gamma\lVert\mathbf x - \mathbf x'\rVert^2)$, which lets the boundary curve
\citep{scholkopf2002}. Only the training points on or inside the margin, the \emph{support vectors}, define the
solution. Support vector regression (SVR) applies the same idea with an $\varepsilon$-insensitive tube around the
regression function.

\paragraph{Experiment.} Because the SVM works with distances, feature scaling is essential. Without it, the
validation ROC-AUC was \SVMRawAUC, barely better than guessing. With standardisation it was \SVMScaledAUC. A
linear kernel (\SVMLinAUC) was slightly better than the default RBF kernel (\SVMRbfAUC). A small grid then
selected $C = \SVMBestC$ and $\gamma = \SVMBestGamma$ (validation ROC-AUC \SVMBestValAUC;
\cref{fig:algsvm}). The small $C$ means a wide, smooth margin works best. The final model, trained on
\SVMTrainRows orders because training time grows roughly with the square of the number of rows, reaches a
test ROC-AUC of \AlgSVMAUC. It keeps \SVMNSupport of its training rows as support vectors, a sign that the two
classes overlap heavily. For problem~A the \SVRBest is better and reaches an MAE of \AlgSVMMAE units
($R^2 = \AlgSVMRtwo$).

\begin{figure}[htbp]\centering
\includegraphics[width=.6\linewidth]{algo/07_svm_tuning}
\caption{Support vector machine: validation ROC-AUC (2024) for combinations of $C$ and $\gamma$}\label{fig:algsvm}
\end{figure}

% ============================================================================
\section{\texorpdfstring{$k$}{k}-nearest neighbours}
\paragraph{Principle.} KNN stores the training data and predicts for a new case from its $k$ closest neighbours
$N_k(\mathbf x)$ in feature space: by majority vote for classification and by the (optionally distance-weighted)
mean for regression \citep{cover1967}:
\[
  \hat y(\mathbf x) = \frac{1}{k}\sum_{i \in N_k(\mathbf x)} y_i .
\]
There is no training step in the usual sense, but every prediction requires distances to all stored rows.

\paragraph{Experiment.} Scaling matters as much as for the SVM. Without it, the validation ROC-AUC was \KNNRawAUC,
because the distance was almost entirely determined by \texttt{units\_lag\_1}, whose values are in the
thousands. With standardisation it rose to \KNNScaledAUC. The choice of $k$ shows the bias--variance trade-off
directly (\cref{fig:algknn}). With $k=1$ the model is perfect on its own data (ROC-AUC \KNNOneFitAUC) but weak on
2024 (\KNNOneValAUC). The validation score kept improving up to the largest value tried, $k=\KNNBestK$: when the
outcome of a single order is mostly noise, averaging over many neighbours helps. The test ROC-AUC is
\AlgKNNAUC. Inspecting the neighbours of a single order shows the explanation KNN offers: ``orders like this one
were large in \emph{x} out of $k$ cases''. For problem~A, $k=\KNNRegK$ distance-weighted neighbours give an MAE
of \AlgKNNMAE units.

\begin{figure}[htbp]\centering
\includegraphics[width=.6\linewidth]{algo/08_choose_k}
\caption{$k$-nearest neighbours: ROC-AUC on the learned data and on 2024 for different $k$}\label{fig:algknn}
\end{figure}

% ============================================================================
\section{\texorpdfstring{$K$}{K}-means clustering}
\paragraph{Principle.} $K$-means partitions $n$ observations into $K$ clusters with centres $\boldsymbol\mu_j$ by
minimising the within-cluster sum of squares \citep{macqueen1967}:
\[
  \min_{C_1,\dots,C_K} \sum_{j=1}^{K}\ \sum_{\mathbf x_i \in C_j} \lVert \mathbf x_i - \boldsymbol\mu_j\rVert^2 .
\]
Lloyd's algorithm alternates two steps until nothing changes: assign each point to its nearest centre, then move
each centre to the mean of its points. Because the result depends on the starting centres, the algorithm is
restarted several times and the best solution is kept.

\paragraph{Experiment.} The experiment first runs Lloyd's algorithm by hand for $K=3$. It converges in \KMIters
iterations, and the within-cluster sum of squares falls at every step. $K$ was then chosen with the elbow plot
and the silhouette coefficient \citep{rousseeuw1987, kaufman1990}. The silhouette selected $K=\KMBestK$ with a
value of \KMSil, which indicates clear structure. The small cluster contains the \KMSmall bread and breakfast
products, which sell fewer units per month, in larger orders and with a higher share of large orders than the
cakes. For $K \geq 3$ the silhouette drops to about \KMSilFour, so the cakes do not form well-separated
sub-groups on these six features. This complements \cref{ch:models}, where a different feature set (revenue,
growth, seasonality and volatility) and a practical range of three to five segments produced a more detailed,
strategy-oriented segmentation. The fitted model can place a new product into a cluster from its expected
profile, so planners can treat it like similar existing products.

\begin{figure}[htbp]\centering
\includegraphics[width=.55\linewidth]{algo/09_choose_k}\hfill
\includegraphics[width=.43\linewidth]{algo/09_clusters}
\caption{$K$-means: elbow and silhouette for $K = 2,\dots,8$ (left); products and cluster centres on the first two principal components (right)}\label{fig:algkm}
\end{figure}

% ============================================================================
\section{Comparative results}
\Cref{tab:algoreg,tab:algocls} bring the supervised experiments together.

\begin{table}[htbp]\centering\small
\caption{Problem A, monthly units per product: test results 2025--2026 of the algorithm lessons (sorted by MAE; best in bold)}\label{tab:algoreg}
\begin{tabular}{@{}lrrrr@{}}
\toprule
\textbf{Algorithm} & \textbf{MAE} & \textbf{RMSE} & \textbf{MAPE (\%)} & \textbf{$R^2$} \\ \midrule
\textbf{XGBoost} & \textbf{96.9} & \textbf{124.9} & \textbf{13.0} & \textbf{0.699} \\
Random Forest & 97.6 & 125.4 & 12.5 & 0.696 \\
SVM & 99.8 & 125.0 & 14.0 & 0.698 \\
Linear Regression & 100.8 & 126.3 & 14.2 & 0.692 \\
Gradient Boosting & 102.6 & 132.1 & 13.9 & 0.663 \\
KNN & 102.6 & 132.1 & 13.9 & 0.663 \\
Decision Tree & 108.0 & 139.8 & 13.5 & 0.623 \\
\midrule
\emph{Baseline: same as last month} & 139.8 & 176.2 & 18.0 & 0.400 \\
\bottomrule
\end{tabular}
\end{table}

\begin{table}[htbp]\centering\small
\caption{Problem B, large order ($\geq 10$ units): test results 2025--2026 of the algorithm lessons (sorted by ROC-AUC; best in bold)}\label{tab:algocls}
\begin{tabular}{@{}lrrrrr@{}}
\toprule
\textbf{Algorithm} & \textbf{accuracy} & \textbf{precision} & \textbf{recall} & \textbf{F1} & \textbf{ROC-AUC} \\ \midrule
\textbf{Logistic Regression} & \textbf{0.592} & \textbf{0.328} & \textbf{0.765} & \textbf{0.459} & \textbf{0.718} \\
XGBoost & 0.632 & 0.348 & 0.719 & 0.469 & 0.717 \\
Decision Tree & 0.635 & 0.349 & 0.711 & 0.468 & 0.712 \\
Random Forest & 0.632 & 0.345 & 0.696 & 0.461 & 0.712 \\
SVM & 0.715 & 0.396 & 0.494 & 0.439 & 0.712 \\
Gradient Boosting & 0.557 & 0.310 & 0.788 & 0.445 & 0.693 \\
KNN & 0.673 & 0.359 & 0.569 & 0.440 & 0.690 \\
\midrule
\emph{Baseline: always small} & 0.774 & 0.000 & 0.000 & 0.000 & 0.500 \\
\bottomrule
\end{tabular}
\end{table}


For problem A every algorithm beats the baseline. The best, \AlgRegBest, reduces the MAE from \AlgRegBaseMAE to
\AlgRegBestMAE units, an improvement of \AlgRegGain\%. The spread between the top models is small, however, and
linear regression (\AlgLinMAE) is within a few units of the ensembles. Only the single decision tree
(\AlgDTMAE) clearly lags behind. For problem B every algorithm lies in a narrow band of ROC-AUC, from
\AlgClsWorstAUC (\AlgClsWorst) to \AlgClsBestAUC (\AlgClsBest). The algorithms differ mainly in where they place
the trade-off between precision and recall, which is a business choice made through the threshold, not a
difference in predictive power.

These results agree with \cref{ch:models}, which used a richer feature set and a different tuning procedure. In
both chapters, (i) all machine-learning models beat the naive baselines, (ii) the leading regressors are
separated by only a few percent, so their exact ranking is not stable, (iii) linear models are competitive with
complex ensembles once good lag features exist, and (iv) large-order prediction is capped at a ROC-AUC of about
0.7 by the information in the data. Reaching the same conclusions with two independent implementations is a
form of triangulation that increases confidence in the findings \citep{makridakis2018, kuhn2019}.

\begin{table}[htbp]\centering\small
\caption{Practical profile of the nine algorithms, based on the experiments}\label{tab:algoprofile}
\begin{tabularx}{\linewidth}{@{}lcccc>{\raggedright\arraybackslash}X@{}}
\toprule
\textbf{Algorithm} & \textbf{Type} & \textbf{Scaling} & \textbf{Explain.} & \textbf{Cost} & \textbf{Best use in this project} \\ \midrule
Linear regression & R & yes\textsuperscript{a} & high & very low & transparent demand baseline \\
Logistic regression & C & yes & high & very low & large-order probability, odds ratios \\
Decision tree & R/C & no & high & low & business rules, explaining drivers \\
Random forest & R/C & no & medium & medium & robust forecasts with ranges \\
Gradient boosting & R/C & no & medium & medium & accurate forecasts (tuned) \\
XGBoost & R/C & no & medium & low & accurate forecasts, early stopping \\
SVM / SVR & R/C & yes & low & high & linear SVR for demand; small data \\
KNN & R/C & yes & medium & high\textsuperscript{b} & case-based explanations \\
$K$-means & U & yes & high & low & product and market segments \\ \bottomrule
\end{tabularx}
\par\smallskip\footnotesize R = regression, C = classification, U = unsupervised; Explain.\ = explainability.
\textsuperscript{a}Not needed for the fit itself, but needed to compare coefficients.
\textsuperscript{b}At prediction time, because every prediction searches all stored rows.
\end{table}

\section{Lessons for practitioners}
\begin{enumerate}
  \item \textbf{Start with a baseline and a simple model.} A naive rule and a linear model set the bar. Here they
        were within a few percent of the best ensemble, at a fraction of the cost and with full explainability.
  \item \textbf{Tune on validation data, never on the test.} Every setting (depth, number of trees, learning
        rate, $C$, $\gamma$, $k$, threshold) was chosen on 2024. The test score is therefore an honest estimate.
  \item \textbf{Scale for distance-based methods.} SVM and KNN dropped to near-random performance without
        scaling. Tree methods do not need it.
  \item \textbf{Watch the gap between training and validation.} Deep trees, $k=1$ neighbours and boosting without
        early stopping all fit the training data almost perfectly and generalise worse.
  \item \textbf{Choose metrics that match the business.} With imbalanced classes, accuracy rewards a useless
        model. ROC-AUC, recall, precision and a deliberately chosen threshold are more meaningful.
  \item \textbf{When all algorithms agree, the limit is the data.} Seven classifiers with the same ROC-AUC show
        that better large-order prediction needs new information, such as customer history, not a different
        algorithm \citep{kelleher2020, geron2022}.
\end{enumerate}
